\documentclass[11pt,a4paper]{article}

\usepackage{amsmath,amssymb,amsthm}
\usepackage{hyperref}
\usepackage[margin=1in]{geometry}
\usepackage{xcolor}
\usepackage{booktabs}
\usepackage{array}
\usepackage{float}
\usepackage{mathtools}

\newtheorem{theorem}{Theorem}[section]
\newtheorem{conjecture}[theorem]{Conjecture}
\newtheorem{corollary}[theorem]{Corollary}
\newtheorem{lemma}[theorem]{Lemma}
\newtheorem{proposition}[theorem]{Proposition}
\newtheorem{definition}[theorem]{Definition}
\theoremstyle{remark}
\newtheorem{remark}[theorem]{Remark}

\newcommand{\PT}{\mathbb{PT}}
\newcommand{\CP}{\mathbb{CP}}
\newcommand{\C}{\mathbb{C}}
\newcommand{\Z}{\mathbb{Z}}
\newcommand{\R}{\mathbb{R}}
\newcommand{\F}{\mathbb{F}}
\newcommand{\cO}{\mathcal{O}}
\newcommand{\cC}{\mathcal{C}}
\newcommand{\cN}{\mathcal{N}}
\newcommand{\cU}{\mathcal{U}}
\newcommand{\cW}{\mathcal{W}}
\newcommand{\cG}{\mathcal{G}}
\newcommand{\cF}{\mathcal{F}}
\newcommand{\PP}{\mathcal{P}}


\title{Stabilizer Algebra and the $k$-Profile Theorem\\
       for Mermin Pentagrams\\
       {\large (Paper XIV)}}

\author{Zhou-Li Chen \\ co-nlang Research}
\date{June 2026}

\begin{document}

\maketitle

\begin{abstract}
We complete the computational foundations for the Weil representation framework
for Mermin pentagrams established in Papers~XI--XIII.
Our main results are:
(1)~the \emph{$k$-Profile Theorem}: all 12{,}096 Mermin pentagrams fall into
exactly 7 $k$-profile types (sorted $V$-intersection dimensions), and every
profile forces odd total parity, upgrading the 500-sample proposition of
Paper~XIII to a full theorem;
(2)~the \emph{Stabilizer Classification}: the four $G_2(2)$ orbits have
stabilizer subgroups $\mathbb{Z}_2$ (Orbit~1) and $S_3$ (Orbits 2, 3, 4),
showing in particular that no orbit has a cyclic stabilizer of order~6;
(3)~the \emph{5-Minus Exclusion}: pentagrams with $k$-profiles $(3,3,3,0,0)$
or $(7,1,1,1,1)$ never attain the all-five-minus h-pattern, with an algebraic
proof for the $(7,1,1,1,1)$ case;
and (4)~\emph{$\beta_{\mathrm{sum}}$ distribution statistics} across all
12{,}096 pentagrams, refining the 500-sample estimates of Paper~XIII.
A key negative result is that the Orbit~4 parity anomaly
(60.6\% 1-minus vs.\ 65--68\% for Orbits 1--3) is \emph{not} explained by
the abstract stabilizer type, since all three order-6 orbits share the same
$S_3$ structure; the anomaly must reside in the embedding of $S_3$ in $G_2(2)$.
We close with an outlook on the Weil representation framework (Paper~XV),
in which $\beta_{\mathrm{sum}} \equiv 2 \pmod{4}$ will be interpreted as a
2-cocycle constraint in the metaplectic extension of $Sp(6,\mathbb{F}_2)$.
\end{abstract}


\section{Introduction}

Papers~XI through~XIII established the parity theorem for Mermin pentagrams
\cite{PaperXI,PaperXII,PaperXIII}.
For each context $C = \{v_1,v_2,v_3,v_4\}$ in a Mermin pentagram, the sign
\[
  s(C) = (-1)^{\beta(C)/2}, \quad
  \beta(C) = \sum_{j<k} \omega_{\mathrm{int}}(v_j, v_k),
\]
gives $\prod_{i=1}^5 s(C_i) = -1$, i.e., $\beta_{\mathrm{sum}} \equiv 2\pmod 4$,
for all 12{,}096 Mermin pentagrams \cite{PaperXI}.
Paper~XIII showed that the Maslov/Kashiwara index framework does not explain
this constraint, identified a parity anomaly in Orbit~4, and upgraded the
$k$-profile result to a 500-sample proposition \cite{PaperXIII}.

This paper resolves the remaining computational questions:

\begin{enumerate}
  \item \textbf{$k$-Profile Theorem} (\S\ref{sec:kprofile}):
        Full 12{,}096-pentagram verification of 7 $k$-profile types, all forcing
        odd parity.

  \item \textbf{Stabilizer Classification} (\S\ref{sec:stabilizer}):
        Computation of the stabilizer subgroup for each of the four
        $G_2(2)$ orbits; result is $\mathbb{Z}_2$ for Orbit~1 and $S_3$ for
        Orbits~2, 3, 4.

  \item \textbf{Parity Distribution and 5-Minus Exclusion}
        (\S\ref{sec:hpattern}):
        Complete h-pattern distribution across all $k$-profiles;
        algebraic exclusion of the all-five-minus pattern for $k$-profile
        $(7,1,1,1,1)$.

  \item \textbf{$\beta_{\mathrm{sum}}$ Distribution} (\S\ref{sec:betadist}):
        Full-enumeration statistics; range $[-34, 30]$, negative bias
        ($\bar{\beta} = -2.5$).
\end{enumerate}

\noindent
All computations were performed by \texttt{stabilizer\_algebra.py} and
\texttt{beta\_distribution.py} in \cite{PaperXIVsuppl}.


\section{\texorpdfstring{The $k$-Profile Theorem}{The k-Profile Theorem}}\label{sec:kprofile}

\subsection{Setup}

Let $V = \mathbb{F}_2^3 \times \{0\}^3 \subset \mathbb{F}_2^6$ be the
standard Lagrangian subspace (the MASA of ``position'' observables).
For a Mermin pentagram $\mathbf{P}$ with Lagrangians $L_1,\ldots,L_5$, define
the \emph{$k$-profile} as the sorted tuple
\[
  \kappa(\mathbf{P}) = (k_1 \geq k_2 \geq \cdots \geq k_5),
  \quad k_i = |L_i \cap V|.
\]
The four possible $k$-types are $k \in \{0,1,3,7\}$, corresponding to the
four MASA types among the 135 Lagrangians of $W(5,\mathbb{F}_2)$: transverse
($k=0$, 64 Lagrangians), one $V$-point ($k=1$, 56), Fano line in $V$ ($k=3$,
14), and $V$ itself ($k=7$, 1).%
\footnote{Paper~XIII used ascending order $(k_1 \leq \cdots \leq k_5)$.
Descending order is adopted here because $k$ measures proximity to $V$,
so larger $k$ is the more ``special'' Lagrangian.
The correspondence is direct reversal:
e.g., Paper~XIII's $(0,0,0,0,1)$ = this paper's $(1,0,0,0,0)$.}

The \emph{integer symplectic form} is defined on
$\mathbb{F}_2^6 \cong \mathbb{F}_2^3 \times \mathbb{F}_2^3$ by
\[
  \omega_{\mathrm{int}}(u,v) = u_x \cdot v_z - u_z \cdot v_x \;\in\; \mathbb{Z},
\]
where $u = (u_x, u_z)$ with $u_x, u_z \in \{0,1\}^3$ are the
position and momentum components and the product is taken over $\mathbb{Z}$
\cite{PaperXI}.
For each context $C_i = \{v_1,v_2,v_3,v_4\} \subset L_i$, define
$\beta(C_i) = \sum_{j<k} \omega_{\mathrm{int}}(v_j,v_k)$ and
$h_i = \beta(C_i)/2 \bmod 2 \in \{0,1\}$.
Odd parity means $\sum_{i=1}^5 h_i \equiv 1\pmod 2$.

\subsection{The Seven Profiles}

\begin{theorem}[k-Profile Theorem]\label{thm:kprofile}
All 12{,}096 Mermin pentagrams fall into exactly 7 $k$-profile types.
Every profile forces odd total parity: $\sum_{i=1}^5 h_i \equiv 1 \pmod 2$.
Verified computationally for all 12{,}096 pentagrams (0 exceptions).
\end{theorem}

\begin{table}[H]
\centering
\caption{The 7 $k$-profiles and their distribution among all 12{,}096 Mermin
pentagrams. The ``Type'' column uses the F2/F4/B classification from Paper~X
\cite{PaperX} where applicable. All profiles give odd parity ($\Sigma h \bmod 2 = 1$).}
\label{tab:kprofiles}
\begin{tabular}{lrrll}
\toprule
$k$-profile (desc.) & Count & \% & Type (Paper~X) & $\Sigma h \bmod 2$ \\
\midrule
$(3,1,1,0,0)$ & 2{,}688 & 22.2\% & F2 & 1 \\
$(1,1,1,0,0)$ & 3{,}584 & 29.6\% & F4 & 1 \\
$(1,0,0,0,0)$ & 2{,}688 & 22.2\% & B  & 1 \\
$(3,3,3,0,0)$ &    896 &  7.4\% & —  & 1 \\
$(3,0,0,0,0)$ &    896 &  7.4\% & —  & 1 \\
$(1,1,1,1,1)$ &    896 &  7.4\% & —  & 1 \\
$(7,1,1,1,1)$ &    448 &  3.7\% & —  & 1 \\
\midrule
Total         & 12{,}096 & 100\% & & \\
\bottomrule
\end{tabular}
\end{table}

\begin{remark}
The three most common profiles (F2, F4, B) together account for 74\% of all
pentagrams and correspond to the type classification introduced in Paper~X.
The remaining four profiles (totalling 26\%) are ``new'' in the sense that they
do not map cleanly to a single Paper~X type (they mix Lagrangians of different
types within a single pentagram).
\end{remark}

\begin{remark}[$k$-profile is $G_2(2)$-equivariant]
Since $G_2(2)$ preserves the symplectic form and the subspace $V$, the
$k$-value $|L \cap V|$ is preserved under the action, so $\kappa(\mathbf{P})$
is a $G_2(2)$-invariant of the pentagram (as a sorted tuple).
The full 12{,}096 enumeration directly verifies that exactly 7 profiles occur,
independently of any orbit-decomposition argument.
\end{remark}


\section{Stabilizer Classification}\label{sec:stabilizer}

\subsection{Setup}

The group $G_2(2)$ has order 12{,}096 and acts on the set $\PP$ of
12{,}096 Mermin pentagrams.
Paper~X \cite{PaperX} established that this action decomposes into 4 orbits:
\[
  \PP = O_1 \sqcup O_2 \sqcup O_3 \sqcup O_4,
  \quad |O_1| = 6{,}048, \quad |O_2| = |O_3| = |O_4| = 2{,}016.
\]
By the orbit-stabilizer theorem, the stabilizer orders are:
\[
  |\mathrm{Stab}(O_1)| = 12{,}096/6{,}048 = 2, \quad
  |\mathrm{Stab}(O_i)| = 12{,}096/2{,}016 = 6 \text{ for } i=2,3,4.
\]
The two groups of order 6 are $\mathbb{Z}_6$ (abelian, element orders $\{1,2,3,6\}$)
and $S_3$ (non-abelian, element orders $\{1,2,2,2,3,3\}$).
They are distinguished by the number of elements of order~2: one for
$\mathbb{Z}_6$, three for $S_3$.

\subsection{Computation}

For each orbit, we computed the stabilizer of a representative pentagram by
testing all 12{,}096 group elements and retaining those fixing all 5 contexts.
The element-order distribution of each stabilizer was computed directly.

\begin{theorem}[Stabilizer Classification]\label{thm:stabilizer}
The stabilizer subgroup of each $G_2(2)$ orbit on $\PP$ is:
$O_1$ has stabilizer $\mathbb{Z}_2$ (abelian, order~2);
each of $O_2$, $O_3$, $O_4$ has stabilizer $S_3$ (non-abelian, order~6,
with 3~elements of order~2 and 2~elements of order~3).
In particular, no orbit has a cyclic stabilizer of order~6.
See Table~\ref{tab:stabilizers} for the complete element-order data.
\end{theorem}

\begin{table}[H]
\centering
\caption{Stabilizer subgroup structure for the four $G_2(2)$ orbits.
Verified computationally for one representative per orbit \cite{PaperXIVsuppl}.}
\label{tab:stabilizers}
\begin{tabular}{lcccccl}
\toprule
Orbit & $|O_i|$ & $|\mathrm{Stab}|$ & Group & Abelian?
      & $\#$(order 2) & $\#$(order 3) \\
\midrule
$O_1$ & 6{,}048 & 2 & $\mathbb{Z}_2$ & Yes & 1 & 0 \\
$O_2$ & 2{,}016 & 6 & $S_3$          & No  & 3 & 2 \\
$O_3$ & 2{,}016 & 6 & $S_3$          & No  & 3 & 2 \\
$O_4$ & 2{,}016 & 6 & $S_3$          & No  & 3 & 2 \\
\bottomrule
\end{tabular}
\end{table}

\begin{remark}[Geometric interpretation]
The $\mathbb{Z}_2$ stabilizer of $O_1$ corresponds to a single symplectic
reflection fixing the pentagram.
The $S_3$ stabilizers correspond to the symmetry group of a triangle
(3-fold rotation and reflection), reflecting the threefold structure
inherent in the 3-qubit setting.
No orbit has a cyclic $\mathbb{Z}_6$ stabilizer; hexagonal symmetry does not
occur among Mermin pentagrams in $W(5,\mathbb{F}_2)$.
\end{remark}

\begin{remark}[Non-conjugate $S_3$ embeddings and Orbit~4 anomaly]
\label{rem:nonconj}
All three order-6 stabilizers are abstractly isomorphic to $S_3$.
However, since $O_2$, $O_3$, $O_4$ are distinct $G_2(2)$-orbits, their
stabilizer subgroups are pairwise non-conjugate in $G_2(2)$: $G_2(2)$
contains at least three distinct conjugacy classes of $S_3$ subgroups.

This implies that the Orbit~4 parity anomaly---Orbit~4 has 60.6\% 1-minus
pentagrams versus 65--68\% for Orbits~1--3 (Paper~XIII \cite{PaperXIII})---is
\emph{not} explained by the abstract stabilizer type, which is $S_3$ for all
three orbits.
The anomaly must originate in the specific \emph{embedding} of $S_3$ into
$G_2(2)$ (i.e., the conjugacy class of the subgroup) and its action on the
parity function, not the abstract group structure.
Identifying which embedding property causes the anomaly is deferred to
future work.
\end{remark}


\section{Parity Distribution and h-Pattern Structure}\label{sec:hpattern}

\subsection{The Three h-Pattern Types}

For a Mermin pentagram, the sorted tuple $(h_1 \leq h_2 \leq \cdots \leq h_5)$
with $h_i \in \{0,1\}$ and $\sum h_i \equiv 1 \pmod 2$ can only be
$(0,0,0,0,1)$, $(0,0,1,1,1)$, or $(1,1,1,1,1)$---corresponding to exactly
1, 3, or 5 minus-contexts.
This is a direct combinatorial consequence of Theorem~\ref{thm:kprofile}
(odd parity) and does not add independent content.
What \emph{is} non-trivial is the distribution of these three types across
$k$-profiles, and the structural exclusion of the 5-minus type for two
specific $k$-profiles.

\subsection{\texorpdfstring{Per $k$-Profile Distribution}{Per k-Profile Distribution}}

\begin{table}[H]
\centering
\caption{Parity distribution (percentage of 1-minus, 3-minus, 5-minus
pentagrams) within each $k$-profile.
All values are computational, from \cite{PaperXIVsuppl}.
The two bottom rows have no 5-minus pentagrams (see
Proposition~\ref{prop:5minus}).}
\label{tab:hpatterns}
\begin{tabular}{lrrrrl}
\toprule
$k$-profile & Count & 1-minus & 3-minus & 5-minus & h-types \\
\midrule
$(3,1,1,0,0)$ & 2{,}688 & 64.7\% & 34.8\% & 0.4\% & 3 \\
$(1,1,1,0,0)$ & 3{,}584 & 63.5\% & 35.5\% & 1.0\% & 3 \\
$(1,0,0,0,0)$ & 2{,}688 & 58.5\% & 40.2\% & 1.3\% & 3 \\
$(1,1,1,1,1)$ &    896 & 72.3\% & 26.8\% & 0.9\% & 3 \\
$(3,0,0,0,0)$ &    896 & 66.1\% & 32.1\% & 1.8\% & 3 \\
$(3,3,3,0,0)$ &    896 & 78.6\% & 21.4\% & ---  & 2 \\
$(7,1,1,1,1)$ &    448 & 78.6\% & 21.4\% & ---  & 2 \\
\midrule
Global        & 12{,}096 & $\approx$65\% & $\approx$34\% & $\approx$1\% & 3 \\
\bottomrule
\end{tabular}
\end{table}

\begin{remark}[Numerical coincidence of $(3,3,3,0,0)$ and $(7,1,1,1,1)$]
\label{rem:coincidence}
The profiles $(3,3,3,0,0)$ and $(7,1,1,1,1)$ share identical parity
distributions (78.6\% 1-minus, 21.4\% 3-minus, 0\% 5-minus) and their
counts satisfy $896 : 448 = 2 : 1$.
This is unlikely to be coincidental.
A speculative connection: both profiles are precisely the two that exclude
5-minus pentagrams (Proposition~\ref{prop:5minus})---the single $k=7$
Lagrangian $V$ and the three $k=3$ Lagrangians with Fano lines in $V$ each
force $h = 0$ on their respective contexts, placing identical structural
constraints on the parity distribution.
The factor-of-2 ratio in counts may reflect the orbit of $V$ itself under a
$\mathbb{Z}_2$ subgroup of $G_2(2)$ relating the two profile types.
\end{remark}

\begin{remark}[Orbit~4 anomaly is intra-profile, not compositional]
\label{rem:orbit4-intraprofile}
The orbit $\times$ $k$-profile cross-table (Table~\ref{tab:crosstable}) shows
that Orbit~4 has the \emph{smallest} fraction of the low-1-minus profile
$(1,0,0,0,0)$ (16.7\%, vs.\ 22.2\% for $O_1$ and 25.0\% for $O_2, O_3$)
and the \emph{largest} fraction of (3,0,0,0,0) (11.1\%, vs.\ 5.6--8.3\%
for $O_1$--$O_3$).
By a weighted-average calculation using the per-$k$-profile 1-minus rates
from Table~\ref{tab:hpatterns}, the $k$-profile composition of Orbit~4 would
predict a 1-minus rate \emph{higher} than the global 65\%, not lower.
The observed anomaly of 60.6\% therefore cannot be attributed to
$k$-profile composition.
\textbf{The Orbit~4 parity anomaly is intra-profile}: Orbit~4 pentagrams
have a systematically lower 1-minus rate \emph{within each $k$-profile type},
relative to the other orbits.
This strengthens the conclusion of Remark~\ref{rem:nonconj}: the anomaly
must originate in the specific $S_3$ embedding in $G_2(2)$, not in the
global pentagram composition.
\end{remark}

\subsection{5-Minus Exclusion}

\begin{proposition}[5-minus exclusion]\label{prop:5minus}
Pentagrams with $k$-profile $(7,1,1,1,1)$ never attain the all-five-minus
h-pattern $(1,1,1,1,1)$.
Pentagrams with $k$-profile $(3,3,3,0,0)$ never attain the all-five-minus
h-pattern; verified computationally for all 896 such pentagrams.
\end{proposition}

\begin{proof}[Proof for $(7,1,1,1,1)$]
The unique Lagrangian with $k=7$ is $V = \mathbb{F}_2^3\times\{0\}^3$ itself.
For any context $C \subset V$, every vector $v \in V$ has $v_z = 0$, so
$\omega_{\mathrm{int}}(u,v) = u_x \cdot v_z - u_z \cdot v_x = 0$ for all
$u, v \in C$.
Therefore $\beta(C) = \sum_{j<k}\omega_{\mathrm{int}}(v_j,v_k) = 0$,
hence $h(C) = 0$ for the $V$-context.
At most four of the five contexts can have $h=1$; the all-five-minus pattern
is impossible.
\end{proof}

\begin{proof}[Proof for $(3,3,3,0,0)$ (computational)]
All 896 pentagrams with $k$-profile $(3,3,3,0,0)$ were verified by exhaustive
enumeration in \cite{PaperXIVsuppl} (Script \texttt{beta\_distribution.py}).
None attains the all-five-minus h-pattern.
An algebraic explanation is deferred to future work; the phenomenon likely
relates to the Fano-line structure forcing $h(C_i) = 0$ for at least two
of the three $k=3$ contexts.
\end{proof}


\section{\texorpdfstring{$\beta_{\mathrm{sum}}$ Distribution}{beta_sum Distribution}}\label{sec:betadist}

\subsection{Global Statistics}

\begin{table}[H]
\centering
\caption{Global $\beta_{\mathrm{sum}}$ statistics over all 12{,}096 Mermin
pentagrams.
Paper~XIII estimated the range from a 500-pentagram sample; the full
enumeration reveals a wider true range.}
\label{tab:betadist}
\begin{tabular}{ll}
\toprule
Quantity & Value \\
\midrule
$\beta_{\mathrm{sum}} \bmod 4$ & $2$ (100\%, all 12{,}096) \\
Range of $\beta_{\mathrm{sum}}$ & $[-34,\; 30]$ \\
Range of $\beta_{\mathrm{sum}}/2$ & $[-17,\; 15]$ \\
Mean $\bar{\beta}_{\mathrm{sum}}$ & $-2.5$ \\
Median & $-2$ \\
Most frequent values & $-2$ (22.3\%), $-6$ (18.6\%), $2$ (18.0\%) \\
\bottomrule
\end{tabular}
\end{table}

\begin{remark}[Range from full enumeration]
The 500-pentagram sample in Paper~XIII \cite{PaperXIII} captured
$\beta_{\mathrm{sum}}/2 \in \{-13,\ldots,+9\}$.
The full 12{,}096 enumeration reveals the true range is $[-17,\;15]$ for
$\beta_{\mathrm{sum}}/2$ (equivalently, $[-34,\;30]$ for $\beta_{\mathrm{sum}}$).
The discrepancy reflects the low frequency of extreme values, which are
unlikely to appear in a 500-sample.
The Kashiwara triple index range $\{-2,\ldots,+2\}$ (Paper~XIII) remains far
smaller than the $\beta_{\mathrm{sum}}/2$ range $[-17,15]$, confirming the
range-mismatch argument for refuting the Kashiwara--$\beta$ identification.
\end{remark}

\begin{remark}[Negative bias]
The mean $\bar{\beta}_{\mathrm{sum}} = -2.5$ and median $= -2$ indicate a
systematic negative bias.
Since $\beta_{\mathrm{sum}} \equiv 2\pmod 4$ and the values $-2, -6, 2$
account for $\approx 59\%$ of all pentagrams, the distribution is concentrated
near $0$ but skewed toward negative values.
This asymmetry is a global property of the standard symplectic form
$\omega_{\mathrm{int}}$ and the choice of $V$; it is not preserved under the
$G_2(2)$ action (which permutes pentagrams but may shift $\beta_{\mathrm{sum}}$).
\end{remark}


\section{Weil Representation Outlook (Paper~XV Preview)}\label{sec:weil}

The results of Papers~XI--XIV collectively establish that the parity theorem
$\beta_{\mathrm{sum}} \equiv 2 \pmod 4$ is:
\begin{itemize}
  \item \emph{computationally universal} (all 12{,}096 pentagrams, Papers~XI, XIV),
  \item \emph{ordering-invariant as a mod-4 quantity} (Paper~XIII),
  \item \emph{holistic}: not decomposable into independent per-context contributions
        (Paper~XIII),
  \item \emph{resistant to known invariants}: Maslov/Kashiwara index,
        $k$-profile, stabilizer type all fail to explain it algebraically.
\end{itemize}

The natural framework for an algebraic proof is the \emph{metaplectic
(Weil) representation} of $Sp(6,\mathbb{F}_2)$ \cite{LionVergne1980,Gerardin}.

\subsection{Metaplectic Extension}

The metaplectic group $\widetilde{Sp}(6,\mathbb{F}_2)$ is the non-trivial
central extension
\[
  1 \;\longrightarrow\; \mathbb{Z}/2 \;\longrightarrow\;
  \widetilde{Sp}(6,\mathbb{F}_2) \;\longrightarrow\;
  Sp(6,\mathbb{F}_2) \;\longrightarrow\; 1,
\]
corresponding to a class in $H^2(Sp(6,\mathbb{F}_2),\,\mathbb{Z}/2)$.
The Weil representation \cite{Gerardin} provides a concrete realization of this
extension over $\mathbb{F}_q$.

\subsection{\texorpdfstring{$\beta$ as a Phase Factor}{beta as a Phase Factor}}

The quantity $\beta(C)/2$ should be interpreted as the \emph{phase factor}
associated to the Lagrangian splitting transition at context $C$.
More precisely, if $\tilde{g}_C \in \widetilde{Sp}(6,\mathbb{F}_2)$ is the
lift of the symplectic transformation associated to the Lagrangian pair
$(L_i, V)$, then $(-1)^{\beta(C)/2}$ is the image of $\tilde{g}_C$ under the
central $\mathbb{Z}/2$.

Under this interpretation, $\beta_{\mathrm{sum}} \equiv 2 \pmod 4$ means that
the product of the five phase factors
\[
  \prod_{i=1}^5 (-1)^{\beta(C_i)/2} = -1
\]
equals the non-trivial central element.
In other words, the five Lagrangian transitions of a Mermin pentagram multiply
to $-1$ in $\widetilde{Sp}(6,\mathbb{F}_2)$---a statement about a 2-cocycle
evaluated on the pentagram loop.

\subsection{Key References for Paper~XV}

An algebraic proof of $\beta_{\mathrm{sum}} \equiv 2\pmod 4$ via this framework
requires:
\begin{enumerate}
  \item Computing $H^2(Sp(6,\mathbb{F}_2),\,\mathbb{Z}/2)$ and identifying the
        Weil cocycle \cite{Gerardin};
  \item Showing that the pentagram $K_5$ loop evaluates to the non-trivial
        class;
  \item Connecting this to the $G_2(2)$ symmetry breaking pattern.
\end{enumerate}
The relevant references are Lion--Vergne \cite{LionVergne1980} (real case,
Maslov cocycle), G\'erardin \cite{Gerardin} (Weil representation over
$\mathbb{F}_q$), and Gurevich--Hadani \cite{GurevichHadani} (canonical
intertwining operators over finite fields).


\section{Open Problems}\label{sec:open}

\begin{enumerate}
  \item \textbf{Algebraic proof of $\beta_{\mathrm{sum}} \equiv 2 \pmod 4$}
        via the Weil 2-cocycle (Paper~XV program).

  \item \textbf{Orbit~4 anomaly}: identify the property of the $S_3$
        embedding in $G_2(2)$ that causes the 60.6\% vs.\ 65--68\% parity
        discrepancy.
        The anomaly is \emph{intra-profile} (Remark~\ref{rem:orbit4-intraprofile}):
        Orbit~4 has a lower 1-minus rate within each $k$-profile type,
        ruling out compositional explanations.
        Candidates: conjugacy class of the $S_3$ subgroup, action on the
        $\beta_{\mathrm{sum}}$ fiber, or cocycle evaluation on the pentagram loop.

  \item \textbf{5-minus exclusion for $(3,3,3,0,0)$}: algebraic proof that
        three Fano-line Lagrangians cannot all contribute $h=1$.

  \item \textbf{$k$-profile algebraic proof}: prove each of the 7 profiles
        forces odd $\Sigma h$ without exhaustive enumeration.
        The $(7,1,1,1,1)$ case is done (Proposition~\ref{prop:5minus});
        the remaining six cases require analysis of $\omega_{\mathrm{int}}$
        on non-V Lagrangians.

  \item \textbf{Class~II geometric characterization}: find a
        $G_2(2)$-equivariant invariant distinguishing Class~I
        ($O_1 \cup O_2$) from Class~II ($O_3 \cup O_4$), beyond the
        F2/F4-type fractions of Paper~X.
        The $O_2$/$O_3$ paradox (identical parity and $q(T)$ distributions
        despite different classes) suggests the invariant must be computed
        at the level of $V$-projections, not Lagrangians.
\end{enumerate}


\appendix
\setcounter{table}{5}
\renewcommand{\thetable}{A\arabic{table}}
\renewcommand{\theHtable}{A\arabic{table}}
\section{Rigor Self-Assessment}

\begin{table}[H]
\centering
\caption{Rigor classification and verification scope for all main results.}
\label{tab:rigor}
\begin{tabular}{p{6cm}ll}
\toprule
Result & Type & Scope \\
\midrule
Thm~\ref{thm:kprofile} (k-profile, 7 types, odd parity)
  & Computational & All 12{,}096 \cite{PaperXIVsuppl} \\
Thm~\ref{thm:stabilizer} ($O_1 = \mathbb{Z}_2$, $O_{2,3,4} = S_3$)
  & Computational & One rep.\ per orbit \\
Rem~\ref{rem:nonconj} (non-conjugate $S_3$ embeddings)
  & Algebraic & Orbit-stabilizer theorem \\
Prop~\ref{prop:5minus} (5-minus excl., $(7,1,1,1,1)$)
  & Algebraic & Exact proof \\
Prop~\ref{prop:5minus} (5-minus excl., $(3,3,3,0,0)$)
  & Computational & All 896 \cite{PaperXIVsuppl} \\
Table~\ref{tab:hpatterns} (parity distribution per $k$-profile)
  & Computational & All 12{,}096 \cite{PaperXIVsuppl} \\
Table~\ref{tab:betadist} ($\beta_{\mathrm{sum}}$ statistics)
  & Computational & All 12{,}096 \cite{PaperXIVsuppl} \\
Table~\ref{tab:crosstable} (orbit $\times$ $k$-profile cross-table)
  & Computational & All 12{,}096 \cite{PaperXIVsuppl} \\
\bottomrule
\end{tabular}
\end{table}

\section{Orbit $\times$ $k$-Profile Cross-Table}

\begin{table}[H]
\centering
\caption{Number of pentagrams in each orbit with each $k$-profile.
All counts are multiples of 56 (= $|\mathcal{P}|/216 = 2^3 \cdot 7$),
reflecting the group-theoretic structure of $G_2(2)$.
Columns are sorted by total count (descending).}
\label{tab:crosstable}
\setlength{\tabcolsep}{4pt}
\begin{tabular}{lrrrrrrr}
\toprule
Orbit & $(1,1,1,0,0)$ & $(1,0,0,0,0)$ & $(3,1,1,0,0)$ & $(3,3,3,0,0)$
      & $(1,1,1,1,1)$ & $(3,0,0,0,0)$ & $(7,1,1,1,1)$ \\
\midrule
$O_1$ (6{,}048) & 1{,}848 & 1{,}344 & 1{,}344 & 504 & 504 & 336 & 168 \\
$O_2$ (2{,}016) &   616 &   504 &   336 & 168 & 112 & 168 & 112 \\
$O_3$ (2{,}016) &   504 &   504 &   504 & 112 & 112 & 168 & 112 \\
$O_4$ (2{,}016) &   616 &   336 &   504 & 112 & 168 & 224 &  56 \\
\midrule
Total & 3{,}584 & 2{,}688 & 2{,}688 & 896 & 896 & 896 & 448 \\
\bottomrule
\end{tabular}
\end{table}

\begin{table}[H]
\centering
\caption{Percentage of each $k$-profile within each orbit
(row percentages, summing to 100\% per orbit).}
\label{tab:crosstable-pct}
\setlength{\tabcolsep}{4pt}
\begin{tabular}{lrrrrrrr}
\toprule
Orbit & $(1,1,1,0,0)$ & $(1,0,0,0,0)$ & $(3,1,1,0,0)$ & $(3,3,3,0,0)$
      & $(1,1,1,1,1)$ & $(3,0,0,0,0)$ & $(7,1,1,1,1)$ \\
\midrule
$O_1$ & 30.6\% & 22.2\% & 22.2\% & 8.3\% & 8.3\% & 5.6\% & 2.8\% \\
$O_2$ & 30.6\% & 25.0\% & 16.7\% & 8.3\% & 5.6\% & 8.3\% & 5.6\% \\
$O_3$ & 25.0\% & 25.0\% & 25.0\% & 5.6\% & 5.6\% & 8.3\% & 5.6\% \\
$O_4$ & 30.6\% & 16.7\% & 25.0\% & 5.6\% & 8.3\% & 11.1\% & 2.8\% \\
\bottomrule
\end{tabular}
\end{table}

\begin{remark}[$O_3$ uniformity and the 56-multiple structure]
Orbit~$O_3$ is the only orbit with equal fractions of the three most common
profiles: $(1,1,1,0,0) = (1,0,0,0,0) = (3,1,1,0,0) = 25.0\%$
(i.e., exactly 504 pentagrams each).
This three-fold balance may reflect an additional symmetry of $O_3$ not
captured by the stabilizer type alone.

All 28 entries in Table~\ref{tab:crosstable} are multiples of 56.
Since $12{,}096 = 216 \times 56$ and the smallest non-zero entry is 56
(the $(7,1,1,1,1)$ count in $O_4$), the natural ``atom'' of this
distribution is $56 = 2^3 \cdot 7$.
This divisibility pattern is a direct consequence of the orbit structure of
$G_2(2)$ and will be relevant to the cocycle computation in Paper~XV.
\end{remark}


\begin{thebibliography}{99}

\bibitem[PaperIX]{PaperIX}
Z.-L. Chen,
``The 3-Qubit Obstruction Ladder: Hyperbolic Embedding,
Mermin Pentagram, and Two Klein Quartic Bridges,''
\textit{Zenodo, DOI: 10.5281/zenodo.20465623} (2026).

\bibitem[PaperX]{PaperX}
Z.-L. Chen,
``The Equiangular Characterization of Mermin Pentagrams:
Symplectic Embedding, 10-Ray Cap, and the Failure of the
Collinearity--Obstruction Dictionary,''
\textit{Zenodo, DOI: 10.5281/zenodo.20476659} (2026).

\bibitem[PaperXI]{PaperXI}
Z.-L. Chen,
``Quadratic Refinement and the Parity of Mermin Pentagrams:
The $\beta$-Formula, Geometric Decomposition, and the
Kochen--Specker Obstruction from Equiangular Geometry,''
\textit{Zenodo, DOI: 10.5281/zenodo.20482595} (2026).

\bibitem[PaperXII]{PaperXII}
Z.-L. Chen,
``The T-Vector Theorem: Symplectic Characteristic Elements
and the Kochen--Specker Obstruction,''
\textit{Zenodo, DOI: 10.5281/zenodo.20490118} (2026).

\bibitem[PaperXIII]{PaperXIII}
Z.-L. Chen,
``The Maslov Index and the Kochen--Specker Obstruction:
Negative Results, the $k$-Profile Theorem, and Orbit~4 Anomaly,''
\textit{Zenodo, DOI: 10.5281/zenodo.20496513} (2026).

\bibitem[PaperXIsuppl]{PaperXIsuppl}
Z.-L. Chen,
``Supplementary computational scripts for Paper~XI,''
\textit{Zenodo, DOI: 10.5281/zenodo.20482283} (2026).

\bibitem[PaperXIIIsuppl]{PaperXIIIsuppl}
Z.-L. Chen,
``Supplementary computational scripts for Paper~XIII,''
\textit{Zenodo, DOI: 10.5281/zenodo.20495857} (2026).

\bibitem[PaperXIVsuppl]{PaperXIVsuppl}
Z.-L. Chen,
``Supplementary computational scripts for Paper~XIV,''
\textit{Zenodo, DOI: 10.5281/zenodo.20501961} (2026).

\bibitem[LionVergne1980]{LionVergne1980}
G. Lion and M. Vergne,
``The Weil Representation, Maslov Index and Theta Series,''
\textit{Birkh\"auser, Boston} (1980).

\bibitem[Gerardin]{Gerardin}
P. G\'erardin,
``Weil representations associated to finite fields,''
\textit{J. Algebra} \textbf{46} (1977), 54--101.

\bibitem[GurevichHadani]{GurevichHadani}
S. Gurevich and R. Hadani,
``Quantization of symplectic vector spaces over finite fields,''
\textit{J. Symplectic Geom.} \textbf{9} (2011), 1--12.

\end{thebibliography}

\end{document}
